\chapter{Forward Curves, Storage and Convenience Yield}\label{ch:m3:forward-curves-storage-and-convenience-yield}

From January 1985 to April 2024 the nearby WTI futures price rose from \$25.92 to \$86.91 a
barrel, a gain of 3.1\% a year. An investor who held the futures over the same years, rolling
each month from the expiring contract to the next, earned less on the oil itself: the roll
cost 0.7\% a year, because the market spent more than half its days in contango and paid
buyers of the curve less than the spot price rose. The difference between what a commodity's
price does and what a futures position in it earns is the shape of its \omterm{def:m3:forward-curves-storage-and-convenience-yield:curve}{forward curve}. This
chapter explains that shape through the economics of storage, measures the \omterm{def:m3:forward-curves-storage-and-convenience-yield:storage}{convenience yield}
the market implies, shows why near contracts are more volatile than far ones, and follows
the money that tracks commodity indices through their monthly roll.

\section{The forward curve of a commodity}

\begin{definition}[Forward curve]\label{def:m3:forward-curves-storage-and-convenience-yield:curve}
The \emph{forward curve}\index{forward curve} of a commodity at a date is the set of prices, on
that date, of futures or forwards on the commodity for successive delivery dates, seen as a function
of the time to delivery.
\end{definition}

A commodity curve differs from a rate curve (One Quant Book 2, chapter 9) in what ties its points
together: not interest alone, but storage. A barrel for delivery in a year is a barrel today plus a
year in a tank, or a barrel produced later. The curve is in contango or in backwardation (One Quant
Book 1, chapter 21) according to which of those two ways of getting the barrel is cheaper.

\begin{omfigure}
\centering
\begin{tikzpicture}[font=\small]
\begin{axis}[width=11cm,height=4.8cm,
  xlabel={time to delivery},ylabel={price (stylised)},
  tick label style={font=\small},label style={font=\small},
  xmin=0,xmax=12,ymin=60,ymax=100,xtick=\empty,ytick=\empty,
  legend columns=3,legend style={font=\footnotesize,at={(0.5,-0.2)},anchor=north,draw=none,fill=none,/tikz/every even column/.append style={column sep=8pt}}]
\addplot[omThm,thick,dashed,domain=0:12,samples=40] {80*exp(0.06*x/12*1.8)};
\addplot[omDef,very thick,domain=0:12,samples=40] {80*exp(0.03*x/12*1.8)};
\addplot[omProp,very thick,domain=0:12,samples=40] {80*exp(-0.12*x/12*1.8)};
\legend{full carry: $r+u$,contango,backwardation}
\end{axis}
\end{tikzpicture}

\omcaption{\omterm{def:m3:forward-curves-storage-and-convenience-yield:curve}{Forward curves}, schematic. Storage caps how steep contango can be (\omterm{def:m3:forward-curves-storage-and-convenience-yield:carry}{full carry}, the cost
of financing and storing), while nothing caps backwardation, which is set by how much the market
values having the commodity now. Stylised, no data.}\label{fig:m3:forward-curves-storage-and-convenience-yield:shapes}
\end{omfigure}

\section{The theory of storage and the convenience yield}

\begin{definition}[Cash-and-carry trade, full carry]\label{def:m3:forward-curves-storage-and-convenience-yield:carry}
A \emph{cash-and-carry trade}\index{cash-and-carry trade} buys a commodity (or a near future) and
sells a later future against it, financing and storing the commodity in between. A curve is at
\emph{full carry}\index{full carry} when the later price exceeds the near one by exactly the cost of
financing and storage, so that the trade earns nothing.
\end{definition}

The \omterm{def:m3:forward-curves-storage-and-convenience-yield:carry}{cash-and-carry trade} bounds contango from above whenever storage is available: $F_{t,T} \le
S_t e^{(r+u)\tau}$, with $u$ the proportional storage cost (the storage floor of
\cref{prop:m3:crude-oil:floor} read from the other side). No trade bounds backwardation: to profit
from a far price below the near one, one would have to sell the commodity now and buy it back later,
which only a holder of inventory can do, and a holder that needs the commodity will not.

\begin{definition}[Theory of storage, convenience yield]\label{def:m3:forward-curves-storage-and-convenience-yield:storage}
The \emph{theory of storage}\index{theory of storage} (Kaldor, Working, Brennan) explains the
futures--spot spread by the cost of storage and a benefit of holding inventory. That benefit, per
unit of value and time, is the \emph{convenience yield}\index{convenience yield} $y_c$: the value of
being able to use the commodity now, avoiding a stock-out, keeping a plant running. With it,
\[
F_{t,T} = S_t\,e^{(r + u - y_c)\,\tau}.
\]
\end{definition}

\begin{proposition}[Convenience yield falls with inventory]\label{prop:m3:forward-curves-storage-and-convenience-yield:inventory}
Between two futures $\Delta\tau$ apart, the net \omterm{def:m3:forward-curves-storage-and-convenience-yield:storage}{convenience yield} implied by the market is
$y_c - u = r - \ln(F_2/F_1)/\Delta\tau$. When inventories are ample, $y_c$ is near zero and the curve
is near \omterm{def:m3:forward-curves-storage-and-convenience-yield:carry}{full carry}; when they are scarce, $y_c$ is large and the curve backwardated.
\end{proposition}

\begin{proof}
Take logs of $F_2/F_1 = e^{(r + u - y_c)\Delta\tau}$. The second statement is the theory's economic
content: the marginal barrel in store is worth more to hold when fewer are held.
\end{proof}

\begin{example}[Reading a spread]\label{ex:m3:forward-curves-storage-and-convenience-yield:read}
The front WTI contract is at \$80.00 and the second at \$79.40, a month apart, with the rate at 4\%.
The net \omterm{def:m3:forward-curves-storage-and-convenience-yield:storage}{convenience yield} is $0.04 - 12\ln(79.40/80.00) = 13.0\%$ a year: the market pays about 13\%
a year above financing to hold oil now. With the second at \$80.60, it is $-5.0\%$: storage is being
paid for.
\end{example}

\begin{omfigure}
\begin{tikzpicture}
\begin{axis}[width=11cm,height=4.8cm,
  ylabel={$y_c-u$, \% a year},
  tick label style={font=\small},label style={font=\small},
  xmin=1985,xmax=2024.4,ymin=-100,ymax=100,ytick={-50,0,50,100},grid=major,grid style={gray!25},
  xtick={1985,1990,1995,2000,2005,2010,2015,2020},xticklabel style={/pgf/number format/1000 sep={}},
  restrict y to domain=-100:100,unbounded coords=jump,
  table/col sep=comma]
\addplot[omDef,thin] table[x=t,y=ycu]{figdata/markets-3/10-forward-curves-storage-and-convenience-yield/convenience.csv};
\draw[gray] (axis cs:1985,0) -- (axis cs:2024.4,0);
\end{axis}
\end{tikzpicture}

\omcaption{Net \omterm{def:m3:forward-curves-storage-and-convenience-yield:storage}{convenience yield} of WTI implied by contracts 1 and 2 and the 3-month Treasury bill,
monthly means of daily values, January 1985 to March 2024 (months below $-100\%$, such as April 2020,
are off the chart). The median month is close to zero; the curve was in backwardation on 45\% of days.
Data: EIA; FRED DTB3.}\label{fig:m3:forward-curves-storage-and-convenience-yield:cy}
\end{omfigure}

When the cash-and-carry pays for more than land tanks, it pays for ships.

\begin{definition}[Floating storage]\label{def:m3:forward-curves-storage-and-convenience-yield:floating}
\emph{Floating storage}\index{floating storage} is oil (or another liquid) held in tankers at anchor
as inventory rather than shipped, which pays when the contango over the months of the charter exceeds
the charter cost plus financing.
\end{definition}

In the contango of 2020 it did: the EIA, citing ClipperData, reports global crude \omterm{def:m3:forward-curves-storage-and-convenience-yield:floating}{floating storage} rising from
65.9 million barrels on 1 March 2020 to 221.5 million on 9 July.

\section{Roll yield and the decomposition of futures returns}

A futures position has no carry of its own: it costs nothing to enter. What an investor earns by
holding futures and rolling them comes from three sources.

\begin{definition}[Spot return, roll yield, collateral return]\label{def:m3:forward-curves-storage-and-convenience-yield:roll}
For a fully collateralised, rolled long futures position, the \emph{spot return}\index{spot return}
is the change in the nearby futures price; the \emph{roll yield}\index{roll yield} is the rest of the
futures' return, earned as each held contract converges to the nearby price (positive in
backwardation, negative in contango); the \emph{collateral return}\index{collateral return} is the
interest earned on the cash that backs the position. Their sum is the total return.
\end{definition}

\begin{proposition}[Roll yield from the curve]\label{prop:m3:forward-curves-storage-and-convenience-yield:rolly}
If the curve's shape is unchanged over a period $\Delta\tau$, a position in a contract $\Delta\tau$
further out than the nearby earns $\ln(F_{\mathrm{near}}/F_{\mathrm{far}})/\Delta\tau$ a year above the
\omterm{def:m3:forward-curves-storage-and-convenience-yield:roll}{spot return}: the \omterm{def:m3:forward-curves-storage-and-convenience-yield:roll}{roll yield} equals the net \omterm{def:m3:forward-curves-storage-and-convenience-yield:storage}{convenience yield} less the rate,
$y_c - u - r$.
\end{proposition}

\begin{proof}
With an unchanged shape the held contract ends where the nearer one started: its log return is the
spot's plus $\ln(F_{\mathrm{near}}/F_{\mathrm{far}})$. Substitute the carry relation.
\end{proof}

Measured on the WTI data, a position that rolls monthly (as indices do, in the window described
below) earned a \omterm{def:m3:forward-curves-storage-and-convenience-yield:roll}{spot return} of 3.08\% a year from January 1985 to April 2024, a \omterm{def:m3:forward-curves-storage-and-convenience-yield:roll}{roll yield} of $-0.71\%$
and a \omterm{def:m3:forward-curves-storage-and-convenience-yield:roll}{collateral return} of 3.18\%: a total of 5.55\% a year
(\cref{fig:m3:forward-curves-storage-and-convenience-yield:index}).

\begin{omfigure}
\begin{tikzpicture}
\begin{axis}[width=11cm,height=5cm,ymode=log,
  ylabel={growth of 1 (log scale)},
  tick label style={font=\small},label style={font=\small},
  xmin=1985,xmax=2024.4,ymin=0.3,ymax=20,grid=major,grid style={gray!25},
  ytick={0.5,1,2,4,8,16},yticklabels={0.5,1,2,4,8,16},
  xtick={1985,1990,1995,2000,2005,2010,2015,2020},xticklabel style={/pgf/number format/1000 sep={}},
  legend columns=2,legend style={font=\footnotesize,at={(0.5,-0.18)},anchor=north,draw=none,fill=none,/tikz/every even column/.append style={column sep=8pt}},
  table/col sep=comma]
\addplot[omDef,thick] table[x=t,y=spot]{figdata/markets-3/10-forward-curves-storage-and-convenience-yield/index.csv};
\addplot[omThm,thick] table[x=t,y=excess]{figdata/markets-3/10-forward-curves-storage-and-convenience-yield/index.csv};
\legend{nearby futures price,rolled futures position (excess return)}
\end{axis}
\end{tikzpicture}

\omcaption{WTI, January 1985 to April 2024: the nearby price and a long position rolled monthly over
the 5th--9th trading days, both starting at 1. The gap is the cumulative \omterm{def:m3:forward-curves-storage-and-convenience-yield:roll}{roll yield}. Contract
identities follow the expiry rule with a weekend-only calendar, so a few roll dates may be off by a
day. Data: EIA contracts 1--4.}\label{fig:m3:forward-curves-storage-and-convenience-yield:index}
\end{omfigure}

\section{The Samuelson effect}

\begin{definition}[Samuelson effect]\label{def:m3:forward-curves-storage-and-convenience-yield:samuelson}
The \emph{Samuelson effect}\index{Samuelson effect} is the tendency of futures prices to be more
volatile the closer they are to delivery.
\end{definition}

Shocks to supply and demand are absorbed by inventory: a shock today moves the near price most, and
later prices less, because stocks and production have time to respond. If the spot's deviation from
its long-run level mean-reverts at speed $\kappa$ (an Ornstein--Uhlenbeck process, One Quant Book 4,
chapter 4), a future with time $\tau$ to delivery moves by $e^{-\kappa\tau}$ times the spot's shock,
and its volatility falls with $\tau$. In the WTI data the annualised volatility of daily changes falls
from 40.9\% for the first contract to 38.3\%, 35.3\% and 33.6\% for the second, third and fourth
(\cref{fig:m3:forward-curves-storage-and-convenience-yield:samuelson}). One Quant Book 6, chapter 16,
builds factor models of the whole curve on this effect.

\begin{omfigure}
\begin{tikzpicture}
\begin{axis}[width=8cm,height=4.4cm,ybar,bar width=18pt,
  xlabel={contract (months to expiry, roughly)},ylabel={annualised vol, \%},
  tick label style={font=\small},label style={font=\small},
  xmin=0.4,xmax=4.6,ymin=0,ymax=48,xtick={1,2,3,4},grid=major,grid style={gray!25},
  nodes near coords,every node near coord/.append style={font=\footnotesize},
  table/col sep=comma]
\addplot[fill=omDef!60,draw=omDef] table[x=contract,y=vol]{figdata/markets-3/10-forward-curves-storage-and-convenience-yield/samuelson.csv};
\end{axis}
\end{tikzpicture}

\omcaption{Annualised volatility of daily log changes of WTI contracts 1 to 4, January 1985 to April
2024, excluding the day each contract rolls and days with non-positive prices. Data: EIA.}\label{fig:m3:forward-curves-storage-and-convenience-yield:samuelson}
\end{omfigure}

\section{Commodity indices and the footprint of their roll}

\begin{definition}[Commodity index, roll window]\label{def:m3:forward-curves-storage-and-convenience-yield:index}
A \emph{commodity index}\index{commodity index} is a published rule for a basket of commodity futures
positions (which contracts, in what weights, rolled when), whose total or excess return investors track
through swaps, funds and notes. Its \emph{roll window}\index{roll window} is the set of days, fixed by
the rule, on which the index sells the expiring contracts and buys the next ones, in equal parts each
day.
\end{definition}

\begin{dated}{2026-09}{Two index roll windows}\label{dat:m3:forward-curves-storage-and-convenience-yield:windows}
The S\&P GSCI rolls over its 5th to 9th business days of each month, 20\% a day; the Bloomberg
Commodity Index over its 6th to 10th business days, also 20\% a day.
\end{dated}

A published window means that everyone knows, weeks ahead, that large index money will sell the near
contract and buy the next one on five known days. Traders can put the same spread on before the window
and take it off during it, selling to the index what it must buy. Mou (2010) measured the price impact
of the Goldman roll from 2000 to March 2010, found strategies front-running it with Sharpe ratios as high
as 4.39, and estimated that index investors forwent 3.6\% of annual return. Several index providers have
since offered variants that roll on other days or choose the contract to hold by the curve's shape.

\begin{proposition}[The cost of a predictable roll]\label{prop:m3:forward-curves-storage-and-convenience-yield:rollcost}
Suppose an index rolls $N$ barrels over $D$ days in equal parts and each day's trade moves the spread
against it by $k$ dollars per million barrels, the move persisting until the window ends. Executing
after its own impact each day, the index pays on average $k\,(N/D)\,(D + 1)/2$ per barrel more than
the spread before the window.
\end{proposition}

\begin{proof}
On day $j$ the spread has moved by $j\,k\,N/D$; the average of $j$ over $1,\dots,D$ is $(D + 1)/2$.
\end{proof}

With 50 million barrels, five days and \$0.004 per million barrels (illustrative), the spread moves
\$0.04 a day and the roll costs \$0.12 a barrel, \$6 million a month.

\section{Tutorial: implied convenience yield and roll yield from four contracts}\label{tut:m3:forward-curves-storage-and-convenience-yield:roll}

\begin{tutorial}
\textbf{Goal.} From the four nearest WTI contracts and the T-bill rate, compute the net \omterm{def:m3:forward-curves-storage-and-convenience-yield:storage}{convenience
yield}, the Samuelson volatilities, and the return of an index-style rolled position split into its
parts. \textbf{End state:}
\cref{fig:m3:forward-curves-storage-and-convenience-yield:cy,fig:m3:forward-curves-storage-and-convenience-yield:index,fig:m3:forward-curves-storage-and-convenience-yield:samuelson}.
\begin{tutsteps}
\item \textbf{Carry arithmetic.} \omterm{def:m3:forward-curves-storage-and-convenience-yield:storage}{Convenience yield}, \omterm{def:m3:forward-curves-storage-and-convenience-yield:carry}{full carry}, \omterm{def:m3:forward-curves-storage-and-convenience-yield:roll}{roll yield} and the roll schedule.
\omcode{code/firm/commcurve/firm_commcurve.py}{36}{56}{Carry relations and the index roll schedule.}\label{lst:m3:forward-curves-storage-and-convenience-yield:carry}
\item \textbf{The rolled position.} Contract identities from the expiry rule, weights from the window.
\omcode{code/markets-3/10-forward-curves-storage-and-convenience-yield/python/m3_curves.py}{71}{90}{A long WTI position rolled over the 5th--9th trading days.}\label{lst:m3:forward-curves-storage-and-convenience-yield:index}
\item \textbf{Run} \texttt{m3\_curves.samuelson()}, \texttt{decomposition()}, \texttt{roll\_cost()} and
\texttt{fig\_curves.py}.
\end{tutsteps}
\textbf{What to change next.} Roll into the fourth contract instead of the second and compare the \omterm{def:m3:forward-curves-storage-and-convenience-yield:roll}{roll
yield}; roll on the last five days before expiry and see what the negative price of April 2020 does to
the position.
\end{tutorial}

\section{Build: the commodity curve object}\label{bld:m3:forward-curves-storage-and-convenience-yield:commcurve}

\begin{build}
\bfield{Purpose} Every commodity desk of the miniature firm prices, hedges and rolls along a \omterm{def:m3:forward-curves-storage-and-convenience-yield:curve}{forward
curve}: the curve object interpolates it, reads its \omterm{def:m3:forward-curves-storage-and-convenience-yield:storage}{convenience yield}, and schedules and measures rolls.
\bfield{Interface} \texttt{ForwardCurve(taus, prices).price(tau)}, \texttt{.shape()};
\texttt{net\_convenience\_yield(f1, f2, dt, r)}; \texttt{full\_carry\_spread}; \texttt{roll\_yield};
\texttt{roll\_weights(business\_day, start, end)}; \texttt{decompose(excess\_log, spot\_log,
collateral\_log)}.
\bfield{Rules} Prices must be positive and maturities increasing (a negative-price day is refused, not
interpolated); interpolation is linear in log price; rates are continuously compounded.
\bfield{Acceptance tests} \texttt{code/firm/commcurve/tests/}: interpolation and shape; the carry
identities; the 20\% a day schedule; the decomposition.
\bfield{Stretch} Seasonal curves (gas, power) with a seasonal factor per month; the curve from futures
with different expiry calendars (Brent against WTI); the index's roll cost measured on real spreads.
\end{build}

\begin{omsources}
\item N.~Kaldor, ``Speculation and economic stability'', \emph{Review of Economic Studies}, 1939;
H.~Working, ``The theory of price of storage'', \emph{American Economic Review}, 1949; M.~Brennan,
``The supply of storage'', \emph{American Economic Review}, 1958.
\item P.~Samuelson, ``Proof that properly anticipated prices fluctuate randomly'', \emph{Industrial
Management Review}, 1965.
\item G.~Gorton and K.~G.~Rouwenhorst, ``Facts and fantasies about commodity futures'', \emph{Financial
Analysts Journal}, 2006.
\item Y.~Mou, ``Limits to arbitrage and commodity index investment: front-running the Goldman roll'',
2010 (SSRN 1716841).
\item S\&P Dow Jones Indices, \emph{S\&P GSCI Methodology}; Bloomberg, \emph{BCOM} methodology.
\item EIA, NYMEX WTI contracts 1--4; FRED series DTB3.
\item US Energy Information Administration, \emph{This Week in Petroleum}, 21 October 2020.
\end{omsources}

\section{Exercises}

\begin{exercise}[$\star$]\label{exo:m3:forward-curves-storage-and-convenience-yield:1}
Spot is \$80, the rate 4\%, storage 3\% a year and the \omterm{def:m3:forward-curves-storage-and-convenience-yield:storage}{convenience yield} 5\%. Give the one-year
forward.
\end{exercise}

\begin{exercise}[$\star$]\label{exo:m3:forward-curves-storage-and-convenience-yield:2}
With the rate at 4\% and storage at 3\%, what is the largest one-month contango on an \$80 barrel?
\end{exercise}

\begin{exercise}[$\star$]\label{exo:m3:forward-curves-storage-and-convenience-yield:3}
Why can backwardation be arbitrarily steep but contango cannot?
\end{exercise}

\begin{exercise}[$\star\star$]\label{exo:m3:forward-curves-storage-and-convenience-yield:4}
The front contract is at \$70.00 and the second at \$71.20, a month apart; the rate is 5\%. Give the
net \omterm{def:m3:forward-curves-storage-and-convenience-yield:storage}{convenience yield} and say what it means.
\end{exercise}

\begin{exercise}[$\star\star$]\label{exo:m3:forward-curves-storage-and-convenience-yield:5}
A rolled position earned 5.55\% a year and the spot rose 3.08\%. With a \omterm{def:m3:forward-curves-storage-and-convenience-yield:roll}{collateral return} of 3.18\%,
what was the \omterm{def:m3:forward-curves-storage-and-convenience-yield:roll}{roll yield}?
\end{exercise}

\begin{exercise}[$\star\star$]\label{exo:m3:forward-curves-storage-and-convenience-yield:6}
In \cref{prop:m3:forward-curves-storage-and-convenience-yield:rollcost}, what does the roll cost per barrel
if the index rolls in one day instead of five? In ten?
\end{exercise}

\begin{exercise}[$\star\star\star$]\label{exo:m3:forward-curves-storage-and-convenience-yield:7}
\emph{Coding.} With \texttt{samuelson}, give the volatilities of contracts 1 to 4 and the ratio of the
fourth to the first. What mean-reversion speed would an Ornstein--Uhlenbeck spot give for that ratio
over three months?
\end{exercise}

\begin{exercise}[$\star\star\star$]\label{exo:m3:forward-curves-storage-and-convenience-yield:8}
\emph{Find the flaw.} ``Oil rose 3\% a year since 1985, so a commodity fund tracking oil futures earned
3\% a year plus interest.''
\end{exercise}

\section{Problem: Front-Running the Roll}

\begin{problem}[{Weekend problem --- an index, its window and the traders ahead of it}]\label{pb:m3:forward-curves-storage-and-convenience-yield:1}
Index money equivalent to 50 million barrels of WTI rolls each month from the second to the third
contract over five known days. Each million barrels traded in a day moves the spread by \$0.004 until
the window ends (illustrative).

\medskip\noindent\textbf{Part I --- The roll.}
\begin{enumerate}
\item How many barrels does the index roll each day?
\item By how much does each day's trade move the spread?
\item What is the index's average cost per barrel against the pre-window spread?
\item And per month, in dollars?
\item And per year?
\end{enumerate}

\medskip\noindent\textbf{Part II --- The front-runner.}
\begin{enumerate}[resume]
\item What position does a trader take before the window?
\item When does it close it, and what does it earn per barrel?
\item What limits how much it can do?
\item What happens to the index's cost as more traders do this?
\item Why is this not manipulation?
\end{enumerate}

\medskip\noindent\textbf{Part III --- Design.}
\begin{enumerate}[resume]
\item What would rolling in one day cost? In ten?
\item Why do index providers keep the window public?
\item How does an index that rolls to deferred contracts change the \omterm{def:m3:forward-curves-storage-and-convenience-yield:roll}{roll yield}?
\item How did Mou measure the impact, and what did he find?
\item What does the index investor see in its returns?
\end{enumerate}

\medskip\noindent\textbf{Part IV --- Judgement.}
\begin{enumerate}[resume]
\item Is a published roll a cost or a service to the index investor?
\item Why does contango make the roll more expensive to index investors even without front-running?
\item Would you invest in a \omterm{def:m3:forward-curves-storage-and-convenience-yield:index}{commodity index} today? What would you check?
\item State the \emph{named result}: the cost of the predictable roll per barrel and per year.
\item In one sentence: who earns the \omterm{def:m3:forward-curves-storage-and-convenience-yield:roll}{roll yield}?
\end{enumerate}
\end{problem}

\section{Interview questions}

\begin{interviewq}[$\star$ \iqroles{trader}]\label{iq:m3:forward-curves-storage-and-convenience-yield:1}
What is the \omterm{def:m3:forward-curves-storage-and-convenience-yield:storage}{convenience yield}, and when is it high?
\end{interviewq}

\begin{interviewq}[$\star$ \iqroles{trader, researcher}]\label{iq:m3:forward-curves-storage-and-convenience-yield:2}
Decompose the return of a commodity futures index.
\end{interviewq}

\begin{interviewq}[$\star\star$ \iqroles{researcher}]\label{iq:m3:forward-curves-storage-and-convenience-yield:3}
Why are near futures more volatile than far ones, and how would you model it?
\end{interviewq}

\begin{interviewq}[$\star\star$ \iqroles{trader}]\label{iq:m3:forward-curves-storage-and-convenience-yield:4}
The curve is in steep contango and tankers are cheap. What trade do you consider, and what can go wrong?
\end{interviewq}

\begin{interviewq}[$\star\star$ \iqroles{researcher, risk}]\label{iq:m3:forward-curves-storage-and-convenience-yield:5}
How would you test whether an index's \omterm{def:m3:forward-curves-storage-and-convenience-yield:index}{roll window} moves prices?
\end{interviewq}

\begin{interviewq}[$\star\star\star$ \iqroles{developer, researcher}]\label{iq:m3:forward-curves-storage-and-convenience-yield:6}
Build a continuous futures series for backtesting. What choices matter and what are the traps?
\end{interviewq}
